\newpage
\phantomsection
\label{prf:supremum-difference-set}
\LRAProofFor{thm:supremum-difference-set}

\begin{remark*}[Return]
\hyperref[thm:supremum-difference-set]{Return to Theorem}
\end{remark*}

\begin{theorem*}[Supremum of a Difference Set]
Let $A,B\subseteq\mathbb{R}$ be nonempty, with $A$ bounded above and $B$ bounded below. Then
\[
  \sup(A-B)=\sup A-\inf B,
  \qquad
  A-B=\{a-b:a\in A,\ b\in B\}.
\]
\end{theorem*}

\begin{proof}[Professional Standard Proof]
\LRAProofBodyStart
Let
\[
  s_A=\sup A
  \qquad\text{and}\qquad
  i_B=\inf B.
\]
For every $a\in A$ and $b\in B$, the definitions of supremum and infimum give
\[
  a\le s_A
  \qquad\text{and}\qquad
  i_B\le b.
\]
Hence $-b\le -i_B$, and therefore
\[
  a-b\le s_A-i_B.
\]
Thus $s_A-i_B$ is an upper bound for $A-B$.

It remains to show that this upper bound is least. Let $\varepsilon>0$. By the
epsilon characterization of the supremum, there exists $a_\varepsilon\in A$
such that
\[
  s_A-\frac{\varepsilon}{2}<a_\varepsilon.
\]
By the epsilon characterization of the infimum, there exists
$b_\varepsilon\in B$ such that
\[
  b_\varepsilon<i_B+\frac{\varepsilon}{2}.
\]
Subtracting the second estimate from the first gives
\[
  s_A-i_B-\varepsilon
  <
  a_\varepsilon-b_\varepsilon.
\]
Since $a_\varepsilon-b_\varepsilon\in A-B$, every positive epsilon produces an
element of $A-B$ above $s_A-i_B-\varepsilon$. Together with the upper-bound
property already proved, the epsilon characterization of the supremum gives
\[
  \sup(A-B)=s_A-i_B=\sup A-\inf B.
\]
\end{proof}

\begin{proof}[Detailed Learning Proof]
\LRAProofBodyStart
Set
\[
  s_A=\sup A
  \qquad\text{and}\qquad
  i_B=\inf B.
\]
We will prove that $s_A-i_B$ is the supremum of the difference set
\[
  A-B=\{a-b:a\in A,\ b\in B\}.
\]
That means we must prove two things: first, $s_A-i_B$ is an upper bound for
$A-B$; second, elements of $A-B$ get arbitrarily close to that upper bound from
below.

For the upper-bound part, take any element of $A-B$. It has the form $a-b$,
where $a\in A$ and $b\in B$. Since $s_A=\sup A$, every element of $A$ lies at
or below $s_A$:
\[
  a\le s_A.
\]
Since $i_B=\inf B$, every element of $B$ lies at or above $i_B$:
\[
  i_B\le b.
\]
Multiplying the second inequality by $-1$ reverses it:
\[
  -b\le -i_B.
\]
Adding $a\le s_A$ and $-b\le -i_B$ gives
\[
  a-b\le s_A-i_B.
\]
So every element of $A-B$ is at most $s_A-i_B$, and $s_A-i_B$ is an upper bound
for $A-B$.

Now we prove that this upper bound is least. Let $\varepsilon>0$. Because
$s_A$ is the supremum of $A$, there is some $a_\varepsilon\in A$ with
\[
  s_A-\frac{\varepsilon}{2}<a_\varepsilon.
\]
Because $i_B$ is the infimum of $B$, there is some $b_\varepsilon\in B$ with
\[
  b_\varepsilon<i_B+\frac{\varepsilon}{2}.
\]
The first choice puts $a_\varepsilon$ close to the top of $A$; the second puts
$b_\varepsilon$ close to the bottom of $B$. Combining the two estimates gives
\[
  a_\varepsilon-b_\varepsilon
  >
  \left(s_A-\frac{\varepsilon}{2}\right)
  -
  \left(i_B+\frac{\varepsilon}{2}\right)
  =
  s_A-i_B-\varepsilon.
\]
Also $a_\varepsilon-b_\varepsilon\in A-B$ by definition of the difference set.
Thus for every $\varepsilon>0$, the set $A-B$ contains an element greater than
$s_A-i_B-\varepsilon$, while $s_A-i_B$ remains an upper bound. This is exactly
the epsilon characterization that $s_A-i_B$ is the supremum of $A-B$.
Therefore
\[
  \sup(A-B)=\sup A-\inf B.
\]
\end{proof}

\begin{remark*}[Proof structure]
The upper-bound half combines $a\le\sup A$ with $-b\le-\inf B$. The leastness
half splits an arbitrary $\varepsilon$ into two halves, chooses
$a_\varepsilon$ within $\varepsilon/2$ of $\sup A$ from below and
$b_\varepsilon$ within $\varepsilon/2$ of $\inf B$ from above, and subtracts
them to produce an element of $A-B$ within $\varepsilon$ of
$\sup A-\inf B$.
\end{remark*}

\begin{dependencies}
\hyperref[def:set-arithmetic-images]{Set Arithmetic Images};
\hyperref[def:supremum]{Supremum};
\hyperref[def:infimum]{Infimum};
\hyperref[def:epsilon-characterization-of-supremum]{Epsilon Characterization of Supremum};
\hyperref[def:epsilon-characterization-of-infimum]{Epsilon Characterization of Infimum}.
\end{dependencies}

\clearpage
